
\documentclass{article}
\usepackage{colortbl}
\usepackage{makecell}
\usepackage{multirow}
\usepackage{supertabular}

\begin{document}

\newcounter{utterance}

\centering \large Interaction Transcript for game `imagegame', experiment `compact\_grids', episode 8 with reka{-}flash{-}3.1{-}t0.0.
\vspace{24pt}

{ \footnotesize  \setcounter{utterance}{1}
\setlength{\tabcolsep}{0pt}
\begin{supertabular}{c@{$\;$}|p{.15\linewidth}@{}p{.15\linewidth}p{.15\linewidth}p{.15\linewidth}p{.15\linewidth}p{.15\linewidth}}
    \# & $\;$A & \multicolumn{4}{c}{Game Master} & $\;\:$B\\
    \hline

    \theutterance \stepcounter{utterance}  
    & & \multicolumn{4}{p{0.6\linewidth}}{
        \cellcolor[rgb]{0.9,0.9,0.9}{
            \makecell[{{p{\linewidth}}}]{
                \texttt{\tiny{[P1$\langle$GM]}}
                \texttt{Let us do the following task. The goal is to fill an empty grid that looks like this:} \\
\\ 
\texttt{▢ ▢ ▢ ▢ ▢} \\
\texttt{▢ ▢ ▢ ▢ ▢} \\
\texttt{▢ ▢ ▢ ▢ ▢} \\
\texttt{▢ ▢ ▢ ▢ ▢} \\
\texttt{▢ ▢ ▢ ▢ ▢} \\
\\ 
\texttt{A filled grid below is 5 by 5 and can look like this:} \\
\\ 
\texttt{▢ ▢ ▢ ▢ ▢} \\
\texttt{▢ ▢ E ▢ ▢} \\
\texttt{▢ ▢ ▢ ▢ ▢} \\
\texttt{▢ ▢ ▢ ▢ ▢} \\
\texttt{X X X X X} \\
\\ 
\texttt{I want you to describe this grid to me, step by step. You don't need to describe the empty squares, which are denoted with "▢". Only describe the location of letters in the grid. Then you wait for me to say "What is your next command?", and then you continue with the next step. Take the size of the grid into consideration while giving commands. You can give only one command at each turn but it can refer to single or multiple cells, rows, columns or even the whole grid at once. When you have described everything, you say "Command: DONE". Only give the command and omit everything else.} \\
\\ 
\texttt{For the filled grid above, here are the example steps.} \\
\\ 
\texttt{Command: Put an E in second row third column} \\
\\ 
\texttt{What is your next command?} \\
\texttt{Command: Fill the last row with X} \\
\\ 
\texttt{What is your next command?} \\
\texttt{Command: DONE} \\
\\ 
\texttt{Ok. Please do this for the following example, which is a 5 by 5 grid.} \\
\texttt{Z ▢ ▢ ▢ Z} \\
\texttt{▢ Z ▢ Z ▢} \\
\texttt{▢ ▢ Z ▢ ▢} \\
\texttt{▢ Z ▢ Z ▢} \\
\texttt{Z ▢ ▢ ▢ Z} \\
\texttt{What is your next command?} \\
            }
        }
    }
    & \\ \\

    \theutterance \stepcounter{utterance}  
    & \multicolumn{4}{p{0.6\linewidth}}{
        \cellcolor[rgb]{0.9,0.9,0.9}{
            \makecell[{{p{\linewidth}}}]{
                \texttt{\tiny{[P1$\rangle$GM]}}
                \texttt{<reasoning>} \\
\texttt{Let's see. The user provided a 5x5 grid filled with some Z's and empty squares. They want me to describe the steps to fill the grid by placing the Z's, similar to the example they gave earlier.} \\
\\ 
\texttt{First, I need to look at the given grid:} \\
\\ 
\texttt{Row 1: Z ▢ ▢ ▢ Z} \\
\texttt{Row 2: ▢ Z ▢ Z ▢} \\
\texttt{Row 3: ▢ ▢ Z ▢ ▢} \\
\texttt{Row 4: ▢ Z ▢ Z ▢} \\
\texttt{Row 5: Z ▢ ▢ ▢ Z} \\
\\ 
\texttt{So, the Z's are forming a sort of diamond or cross shape. Let me visualize this. The first and last rows have Z's in the first and last columns. The middle row (row 3) has a Z in the center. Then rows 2 and 4 have Z's in the second, third, fourth columns, but offset. Wait, let me check each row:} \\
\\ 
\texttt{Row 1: Z in column 1 and 5, others empty.} \\
\texttt{Row 2: Z in column 2 and 4, others empty.} \\
\texttt{Row 3: Z only in column 3.} \\
\texttt{Row 4: Same as row 2, Z in 2 and 4.} \\
\texttt{Row 5: Same as row 1, Z in 1 and 5.} \\
\\ 
\texttt{So the pattern is that each row alternates between having Z's in certain columns. The Z's form a diamond shape centered in the grid. The middle row (row 3) has a single Z in the center. The rows above and below (rows 2 and 4) have Z's two positions from the center in their respective rows. The first and last rows have Z's at the very ends.} \\
\\ 
\texttt{To describe this step by step, I need to identify the positions of each Z and give commands to place them. The example given in the problem used commands like "Put an E in second row third column" and "Fill the last row with X". So, similar commands here.} \\
\\ 
\texttt{Looking at the Z positions:} \\
\\ 
\texttt{Row 1: columns 1 and 5.} \\
\texttt{Row 2: columns 2 and 4.} \\
\texttt{Row 3: column 3.} \\
\texttt{Row 4: columns 2 and 4.} \\
\texttt{Row 5: columns 1 and 5.} \\
\\ 
\texttt{So, for each row, I can either specify individual cells or find} \\
            }
        }
    }
    & & \\ \\

    \theutterance \stepcounter{utterance}  
    & & & \multicolumn{2}{p{0.3\linewidth}}{
        \cellcolor[rgb]{0.9,0.9,0.9}{
            \makecell[{{p{\linewidth}}}]{
                \texttt{\tiny{[GM$|$GM]}}
                \texttt{Invalid instruction format} \\
            }
        }
    }
    & & \\ \\

\end{supertabular}
}

\end{document}
